\subsubsection{Comparison between JPEG and JP2K (JPEG2000)}
\begin{itemize}
    \item For high bitrates: no lots of differences
    \item For low bitrates: JPEG becomes very bad but still JPEG2000 has a better quality
    \begin{itemize}
        \item Ringing Artifacts
        \item Even comparing 0.2bpp JPEG with 0.1bpp JP2K, this last one is better visually
    \end{itemize}
    \item JP2K allows the reachability of every code rate we want.
\end{itemize}
\subsubsection{Error Robustness in compressed data}
Error robustness is the ability of a coded stream to contain the spatial or temporal effect of corrupted and missing bits. It is commonly evaluated against bit-error or packet-loss probability and usually costs extra rate through markers, resynchronization, redundancy, or independent coding units.
Errors can be introduced in TXing information, or saving it into files.
Compressed streams are vulnerable, can induce error propagation. We need some `marker' sequences of bits to recover errors.
Tradeoff: robustness-rate.
\begin{itemize}
    \item JPEG with $P_e=10^{-4}$: corrupted bit = all the following part of the image corruption
    \item JPEG2000 with $P_e=10^{-4}$: corrupted bit = loss of detail in $\geq 1$ subbands, visually no great loss.
\end{itemize} 
\subsection{Conceptual Maps}
\begin{minipage}{0.6\textwidth}
    \textbf{DWT}
    \begin{itemize}
        \item Signal model
        \item Filter Banks (analysis/synthesis)
        \item 2D-DWT: approximation, detail of the subbands
        \item Energy Concentration (sparsity), intra/inter band correlation
    \end{itemize}    
\end{minipage}
\begin{minipage}{0.35\textwidth}
    \textbf{EZW/JP2K}
    \begin{itemize}
        \item Wavelet Transform
        \item Bitplane Coding
        \item Coding Strategy
        \begin{itemize}
            \item EZW
            \item JP2K
        \end{itemize}
        \item Comparison with JPEG
    \end{itemize}    
\end{minipage}
\section{Learned Image Coding (LIC) / Neural Image Coding (NIC)}
Learned image coding uses neural networks to learn analysis transforms, synthesis transforms, and probability models jointly from training images. Unlike a fixed handcrafted transform, its parameters are optimized end to end for a selected rate--distortion objective, often using VAEs\footnote{Variational Auto-Encoders} or GANs\footnote{Generative Adversarial Networks}.
\subsection{Coding Architecture}
The analysis transform $g_a$ maps pixels to a latent representation designed for quantization and entropy coding. The synthesis transform $g_s$ reconstructs pixels from quantized latents; together they play roles analogous to transform and inverse transform in a conventional codec.
$$\text{original }x \qquad \longrightarrow \qquad \boxed{\stackrel{\text{Analysis}}{z = g_a(x)}} \qquad \longrightarrow \qquad \boxed{\stackrel{\text{Quantization}}{\hat{z} = Q(z)}} \qquad \longrightarrow \qquad \boxed{\stackrel{\text{Synthesis}}{\hat{x} = g_s(\hat{z})}} \qquad \longrightarrow \qquad \hat{x} \text{ decoded}$$
Legacy framework (JPEG): $g_a$ is a Block Based DCT, while $g_s$ is a DWT.\\
Learned Paradigms allow the use of the optimal $g_a, g_s$ by learning directly from data.
\subsubsection{Optimization of the Loss Function}
The loss function is the scalar objective minimized during training. $D$ penalizes reconstruction error, $R$ estimates expected coded length in bits, and $\lambda$ sets how much rate is traded for quality.
It all goes toward the optimization of the loss function:
$$\mathcal{L} = D(x,\hat{x}) + \lambda R(\hat{z})$$
\subsection{NN Recap}
For image processing we use CNNs.
\subsubsection{Gradient Descent and Backpropagation}
$$\theta_{t+1} = \theta_t - \eta\nabla_\theta \mathcal{L}\left(\theta_t\right)$$
\subsubsection{GDN: Generalized Divisive Normalization}
Generalized Divisive Normalization (GDN) is a learned nonlinear normalization that scales each feature using the energy of neighboring feature channels. It helps reduce statistical dependencies and makes latents easier to entropy-code.
$$w_i = \frac{v_i}{\sqrt{\beta_i + \sr{}{j}\gamma_{ij}v_j}}$$
Given a feature $v_i$ the parameters $\beta_i, \gamma_{ij}$ are learned. This allows us to:
\begin{itemize}
    \item Represent Masking effect
    \item Linking to the perceptron
    \item Gaussianization of the data (decorrelation of the signal = use entropy coding)
\end{itemize}
\subsection{JPEG-AI standard}
It has nothing to do with the original JPEG. Standardized as ISO/IEC 29170-2.
\subsubsection{Core Idea}
Going from pixels to 3D-Latent Tensor (sparse signal). JPEG AI uses $\begin{cases}
    160 \text{ Luminance channels (human sensitivity)}\\
    96\text{ Chrominance channels}
\end{cases}$
\subsubsection{Auto-Encoder Framework}
\begin{itemize}
    \item Encoder: Analysis Transform ($g_a$)
    \item Bottleneck: Latents $y$ are quantized to $\hat{y}$
    \item Decoder: Synthesis Transform ($g_s$)
\end{itemize}
\subsubsection{Rate-Distorsion Variable Auto-Encoder}
In a learned codec, rate is the expected number of bits required to entropy-code quantized latents under the learned probability model. It is usually reported as bpp for images; distortion is an expected fidelity loss such as MSE, MS-SSIM loss, or a task-specific metric.
Objective: minimizing $\mathcal{L}$ defined as
$$\mathcal{L} = R + \lambda D \qquad \text{where }\begin{cases}
    \text{Rate }R = D_{KL}\left[q\left(\hat{y}|x\right)\left|\right|p\left(\hat{y}\right)\right]\\
    \text{Distorsion }D = E\left[\rho(x,\hat{x})\right] 
\end{cases}$$
Tradeoff: higher $\lambda$ = higher quality, lower $\lambda$ = high compression. Compressing always in latent space.\\
Problem: Non differentiability for the quantization function, backprop gradient not working properly. We need:
\begin{itemize}
    \item \textbf{Hard} quantization at \textbf{test} time 
    \item \textbf{Soft} quantization at \textbf{training} time $\to$ Additive Uniform Noise
\end{itemize}
Additive uniform noise has the statistical description of the quantization, and makes the backpropagation work.
$$\hat{z} \approx z + u \qquad u \sim \mathcal{U}\left(-0.5, 0.5\right)$$
\subsubsection{Scale Hyperprior: spatial adaptability}
A scale hyperprior is transmitted side information describing local latent statistics, such as coefficient scales. Although the hyperprior consumes bits, it can lower total rate by making the main latent probability model more accurate at each spatial location.
IDEA: set NN to learn statistics of the quantized data, additional side information $\hat{z}$ to predict $p_{\hat{y}}$ $\forall$ location.\\
We change paradigm: since JPEG-AI is a Nonlinear Tranform Coding, we can outperform HEVC and VVC (intra coding) by 50\% in coding efficiency. We need a Hierarchical VAE model.
\subsubsection{Hierarchical VAE}
Multimodality: compression can be adapted do human consumption or machine consumption. JPEG-AI (at 0.5bpp) saves up to $27\%$ compared with VVC. HF are preserved but huge complexity.\\
Complexity is measured in kilo Multiplication-Accumulation-Count over pixel (kMAC/px). There are 3 profiles:
\begin{itemize}
    \item Dec0 $\to 8$ kMAC/px (low-end CPUs)
    \item Dec1 $\to 23$ kMAC/px (also for Mid-range smartphones)
    \item Dec2 $\to 214$ kMAC/px (high-end GPUs)
\end{itemize}
\subsubsection{Conclusions}
\begin{itemize}
    \item E2E Optimization: $\mathcal{L} = R+\lambda D$
    \item Data driven transforms (Non linear way)
    \item Current Challenges:
    \begin{itemize}
        \item Energy Consumption
        \item Cross Platform Determinism
        \item Out of distribution vulnerability (Hallucinations and artifacts)
    \end{itemize}
    \item State of the art: tradeoff
    \begin{itemize}
        \item Pros $\to$ R-D efficiency
        \item Cons $\to$ COMPLEXITY
    \end{itemize}
\end{itemize}
